\chapter{Answers to the questions}
\label{app:answers}

\section*{\Cref{ch:intro}}
\begin{enumerate}
  \item It maps PWM straight to speed with fitted constants and no physical structure. It
        assumes the speed jumps instantly to its steady value, so it cannot represent
        inertia (acceleration and coasting), nor voltage dependence or a dead band.
  \item Its structure is physical, so it extrapolates sensibly to unseen inputs, and its
        parameters can be sanity-checked against physics and literature.
\end{enumerate}

\section*{\Cref{ch:frames}}
\begin{enumerate}
  \item $N$ flips and $v$ flips, so the coefficient is unchanged: $N_v = \qty{-2.0}{\newton\second}$.
  \item They are diagonal (force index and velocity index belong to the same axis), so both
        signs flip together, or neither does.
  \item Drag depends on velocity relative to the body, so it needs the twist in \code{base_link}.
        Comparing tracks needs the pose in \code{odom}.
\end{enumerate}

\section*{\Cref{ch:fossen}}
\begin{enumerate}
  \item $(\vect{\omega}\times\vect{v})_x = q w - r v$, so $m(\dot u + qw - rv) = X$.
  \item $\mtot = 13.5 - 6.36 = \qty{7.14}{\kilogram}$. That is positive, so it is \emph{not}
        rejected, and the model is silently about \qty{64}{\percent} too light. The positivity check only
        catches extreme cases.
  \item Weight and buoyancy are both vertical, so rotation about the vertical axis does not change
        their moment arms. No yaw angle is preferred.
  \item It enters through $K_\phi = M_\theta = B\,\overline{BG}$, in \code{restoring.K_phi} and
        \code{restoring.M_theta}. $\overline{BG}$ itself is not stored. (Since 15 October 2026
        only \code{restoring.M_theta} remains, because roll is not modelled; see \cref{sec:no-roll}.)
\end{enumerate}

\section*{\Cref{ch:symmetry}}
\begin{enumerate}
  \item $X_{vv}v^2$ is even in $v$, so it is allowed. $X_v v$ is odd, so it is forbidden.
  \item $r$ is odd under the reflection and cannot appear in surge for a symmetric vehicle, while $r^2$
        is even. Testing $r$ is still useful as a sanity check: a significant $r$ means the
        vehicle or its thrust is \emph{not} symmetric (e.g.\ a misaligned thruster or tether
        pulling sideways), or the residual contains a controller artefact.
  \item Nonlinear cross-group terms such as $ur$ and $vr$ (type 2). Combined surge--yaw also
        exercises any actuation-level coupling (type 3) between the surge and yaw thrusters.
\end{enumerate}

\section*{\Cref{ch:surge}}
\begin{enumerate}
  \item Set $\dot u = 0$ in \eqref{eq:surge} and solve the quadratic. $\mtot$ multiplies
        $\dot u$, so it disappears.
  \item Quadratic: $1/u$ against $t$, with slope $\Xuu/\mtot$. Linear: $\ln u$ against $t$, with
        slope $-X_u/\mtot$.
  \item They fit identically: the second is the first scaled by $1.1$, which is exactly the
        unobservable scale factor.
\end{enumerate}

\section*{\Cref{ch:actuation}}
\begin{enumerate}
  \item As the battery sags the true thrust falls by up to about \qty{25}{\percent}, but a
        fixed-voltage curve keeps reporting the same force. The fit explains the lower speed for
        the ``same'' force with more drag.
  \item The propeller is optimised for one direction. A reversed thruster driven ``forward''
        runs on the reverse branch of the curve, which mirroring selects and negation does not.
  \item $0.03 \times 20 = \qty{0.6}{\newton\meter}$, positive about $+y$ in FLU, so nose down.
\end{enumerate}

\section*{\Cref{ch:integration}}
\begin{enumerate}
  \item From $2.7\times10^{-9}$ to $1.7\times10^{-10}$, a factor of about 16, which equals
        $2^4$. This is consistent with fourth order.
  \item The ESC itself holds the last pulse width until a new one arrives (ZOH), so holding is
        the physically correct input model. Interpolating would also use a future sample.
  \item If $\taueff$ became shorter than about \qty{4}{\milli\second}, i.e.\ with enormous damping or a
        tiny inertia, such as a wrongly entered coefficient off by orders of magnitude.
\end{enumerate}

\section*{\Cref{ch:maneuvers}}
\begin{enumerate}
  \item One run gives one equation $T = X_u u + \Xuu u^2$ with two unknowns.
  \item $\delta_1 = \ln(0.30/0.12) = 0.92$, so $\zeta = \delta/\sqrt{4\pi^2 + \delta^2} \approx 0.14$.
        $\delta_2 = \ln(0.12/0.05) = 0.88$ is slightly smaller at smaller amplitude, which hints at a
        small quadratic contribution.
  \item With $T = 0$ there is no thrust term in the equation. It still determines only the
        ratios $X_u/\mtot$ and $\Xuu/\mtot$, not absolute drag.
\end{enumerate}

\section*{\Cref{ch:timing}}
\begin{enumerate}
  \item Inertia is read from the first second of a transient, where a few hundred
        milliseconds matter. Drag is read from steady states or slow decays, where they do not.
  \item A thrust step is an edge in acceleration but only the start of a ramp in velocity, and the
        IMU samples about ten times faster.
  \item At rest $a_x = 0$, so $f_x = -g\sin\theta_0 = -9.81\sin\qty{1}{\degree} \approx
        \qty{-0.17}{\meter\per\second\squared}$. Yes, it is consistent.
\end{enumerate}

\section*{\Cref{ch:estimation}}
\begin{enumerate}
  \item From the normal equations, $\vect{\Phi}^\top(\vect{y} - \vect{\Phi}\hat{\vect{\theta}}) = \vect{0}$.
  \item A causal filter delays $u$ (and $\dot u$) relative to $T$, and the delay looks like extra
        inertia (\cref{ch:timing}).
  \item No differentiation of noisy velocity is needed, and the delay can be estimated as a
        parameter. Drag build-up during the step is also modelled instead of ignored.
\end{enumerate}

\section*{\Cref{ch:coupling}}
\begin{enumerate}
  \item The least-squares residual is already orthogonal to the own regressors, so projecting it
        again would change nothing.
  \item $3/\sqrt{149} \approx 0.25$.
  \item A closed-loop controller that reacts to both surge and sway errors, a current or tether
        whose effect changes with heading and drift, or a misaligned thruster.
\end{enumerate}

\section*{\Cref{ch:validation}}
\begin{enumerate}
  \item A free run hides nothing: every coefficient error accumulates. Re-initialisation
        repeatedly removes accumulated error and only tests short-horizon prediction, although
        its error shape tells which coefficient is wrong.
  \item A velocity bias means thrust scale or drag. Steady-state runs at several levels compare
        $u_\mathrm{ss}$ directly, and a coast-down tests drag without any thrust.
  \item The sign check would show a negative correlation between $(p - p_0)$ and EKF surge and
        print \code{INVERTED - consider thrust.sign: -1}. It would not change anything itself.
\end{enumerate}

\section*{\Cref{ch:synthetic}}
\begin{enumerate}
  \item On the ideal bag every deviation from the truth is a bias of the method itself. If that is
        not zero first, the deviations on imperfect bags mix two causes that can no longer be
        separated.
  \item $\code{servo_delay_s} = \qty{-0.12}{\second}$: the scripts subtract the value from the log
        time, and the thrust acts after the log.
  \item Two errors happen to cancel: the smoothing bias (about $+\qty{14}{\percent}$) and the thrust
        scale ($1/1.15$). $X_u$ and $\Xuu$ of the same fit are off, and a different window would
        break the cancellation.
  \item A short window passes more of the differentiation noise into $\dot u$. Noise in a
        regressor attenuates its coefficient toward zero (errors-in-variables), so $\mtot$ is
        underestimated.
\end{enumerate}
